\documentclass[aspectratio=169,11pt]{beamer}
% Shared CMPSC 480 Beamer preamble.
\usepackage[T1]{fontenc}
\usepackage[utf8]{inputenc}
\usepackage{lmodern}
\usepackage{amsmath,amssymb,mathtools}
\usepackage{booktabs,array,tabularx}
\usepackage{xcolor}
\usepackage{listings}
\usepackage{tikz}
\usetikzlibrary{arrows.meta,positioning,shapes.geometric}

% Ignore only negligible TeX box diagnostics that are below visual tolerance.
% Larger layout defects still appear in the build log and in rendered QA.
\vfuzz=3pt
\hbadness=2000

\definecolor{courseink}{HTML}{111111}
\definecolor{coursegray}{HTML}{EDEDED}
\definecolor{courserule}{HTML}{B8BCC4}
\definecolor{courseblue}{HTML}{3D8DFF}
\definecolor{courselightblue}{HTML}{D0EDFA}
\definecolor{coursemuted}{HTML}{5C6470}
\definecolor{coursegreen}{HTML}{0B7A53}
\definecolor{coursered}{HTML}{B3261E}

\usetheme{default}
\usefonttheme{professionalfonts}
\setbeamertemplate{navigation symbols}{}
\setbeamersize{text margin left=0.65cm,text margin right=0.65cm}

\setbeamercolor{normal text}{fg=courseink,bg=white}
\setbeamercolor{frametitle}{fg=courseink,bg=white}
\setbeamercolor{title}{fg=courseink,bg=white}
\setbeamercolor{subtitle}{fg=coursemuted,bg=white}
\setbeamercolor{structure}{fg=courseblue}
\setbeamercolor{alerted text}{fg=coursered}
\setbeamercolor{block title}{fg=courseink,bg=courselightblue}
\setbeamercolor{block body}{fg=courseink,bg=coursegray}
\setbeamercolor{example text}{fg=coursegreen}

\setbeamerfont{title}{size=\fontsize{25}{29}\selectfont,series=\bfseries}
\setbeamerfont{subtitle}{size=\fontsize{13}{16}\selectfont}
\setbeamerfont{frametitle}{size=\fontsize{19}{22}\selectfont,series=\bfseries}
\setbeamerfont{normal text}{size=\fontsize{11}{14}\selectfont}
\setbeamerfont{footline}{size=\fontsize{7.5}{9}\selectfont}

\setbeamertemplate{frametitle}{%
  \nointerlineskip
  % Use content-driven height so a long assessment title can wrap without
  % being clipped by a fixed-height title bar.
  \begin{beamercolorbox}[wd=\paperwidth,sep=0.17cm,leftskip=0.48cm,rightskip=0.48cm]{frametitle}
    \insertframetitle
  \end{beamercolorbox}
  \vspace{-0.08cm}
  {\color{courserule}\rule{\textwidth}{0.35pt}}
}

\setbeamertemplate{footline}{%
  \leavevmode
  \begin{beamercolorbox}[wd=\paperwidth,ht=2.6ex,dp=1.1ex,leftskip=.65cm,rightskip=.65cm]{author in head/foot}
    \color{coursemuted}\insertshorttitle\hfill\insertframenumber/\inserttotalframenumber
  \end{beamercolorbox}
}

\setbeamertemplate{itemize item}{\color{courseblue}\small$\blacktriangleright$}
\setbeamertemplate{itemize subitem}{\color{courseblue}\scriptsize$\bullet$}
\setbeamertemplate{enumerate item}{\color{courseblue}\insertenumlabel.}

\lstdefinestyle{coursepython}{
  language=Python,
  basicstyle=\ttfamily\fontsize{8.5}{10}\selectfont,
  keywordstyle=\color{courseblue}\bfseries,
  commentstyle=\color{coursemuted}\itshape,
  stringstyle=\color{coursegreen},
  showstringspaces=false,
  columns=fullflexible,
  keepspaces=true,
  frame=single,
  rulecolor=\color{courserule},
  backgroundcolor=\color{coursegray},
  breaklines=true,
  tabsize=4,
  xleftmargin=0.15cm,
  xrightmargin=0.15cm,
  framexleftmargin=0.15cm,
  framexrightmargin=0.15cm
}
\lstset{style=coursepython}

\newcommand{\points}[1]{\hfill{\color{courseblue}\textbf{[#1 points]}}}
\newcommand{\deliverable}[1]{\textbf{\color{courseblue}#1}}
\newcommand{\code}[1]{\texttt{#1}}
\newcommand{\keyidea}[1]{%
  \begin{beamercolorbox}[rounded=false,sep=0.55em,wd=\linewidth]{block body}
    \textbf{Key idea.} #1
  \end{beamercolorbox}
}
\newcommand{\answerlabel}{\textbf{\color{coursegreen}Answer.}\ }
\newcommand{\gradinglabel}{\textbf{\color{courseblue}Grading guidance.}\ }

\newenvironment{questionblock}[1]
  {\begin{block}{#1}}
  {\end{block}}

\newenvironment{solutionblock}[1]
  {\begin{block}{#1}}
  {\end{block}}

\AtBeginSection[]{
  \begin{frame}
    \vfill
    {\usebeamerfont{title}\color{courseink}\insertsectionhead\par}
    \vspace{0.4cm}
    {\color{courseblue}\rule{0.32\paperwidth}{3pt}}
    \vfill
  \end{frame}
}


% CMPSC480_TOTAL_POINTS: 20
% CMPSC480_ITEM_IDS: 1,2,3,4
\title[Quiz 2 Questions]{Quiz 2: Adversarial and Sequential Decision Making}
\subtitle{Weeks 3-5}
\author{CMPSC 480 - Student Question Deck}
\date{}

\begin{document}


\begin{frame}
  \titlepage
\end{frame}

\begin{frame}{Assessment overview}
\begin{columns}[T,onlytextwidth]
\begin{column}{0.62\textwidth}
\Large Trace adversarial backups, solve a Bellman step, and apply a Q-learning update.

\vspace{0.35cm}
\normalsize
Coverage: Weeks 3-5
\end{column}
\begin{column}{0.33\textwidth}
\begin{block}{Points}20 total\end{block}
\begin{block}{Expected time}35 minutes\end{block}
\begin{block}{Submission}Individual PDF\end{block}
\end{column}
\end{columns}
\end{frame}

\begin{frame}{Learning objectives}
\begin{itemize}
  \item Trace minimax and pruning.
  \item Interpret an MDP.
  \item Perform Bellman and Q-learning updates.
  \item Distinguish planning from learning.
\end{itemize}
\end{frame}

\begin{frame}{Requirements checklist}
\small
\begin{itemize}
  \item Work independently.
  \item Show backed-up values and arithmetic.
  \item State terminal conventions.
  \item No electronic communication.
\end{itemize}
\end{frame}

\begin{frame}{Task 1 - Minimax and alpha-beta (6 points)}
\small
A MAX root has MIN children A and B. A has leaves 3 and 5; B has leaves 2 and 9.

\vspace{0.2cm}
\begin{enumerate}
\item[\textbf{a.}] Compute the minimax value and selected branch. \hfill{\color{courseblue}\textbf{[2 points]}}
\item[\textbf{b.}] Search A first. After B's first leaf 2, can B's second leaf be pruned? \hfill{\color{courseblue}\textbf{[2 points]}}
\item[\textbf{c.}] Does pruning change the backed-up value? Explain. \hfill{\color{courseblue}\textbf{[2 points]}}
\end{enumerate}
\end{frame}

\begin{frame}{Task 2 - MDP backup (6 points)}
\small
From state s, action a reaches terminal +5 with probability 0.8 and state t with probability 0.2. The immediate reward is 0, gamma=0.9, and V(t)=2.

\vspace{0.2cm}
\begin{enumerate}
\item[\textbf{a.}] Write the Bellman action-value expression. \hfill{\color{courseblue}\textbf{[2 points]}}
\item[\textbf{b.}] Compute Q(s,a). \hfill{\color{courseblue}\textbf{[2 points]}}
\item[\textbf{c.}] Explain what value iteration maximizes. \hfill{\color{courseblue}\textbf{[2 points]}}
\end{enumerate}
\end{frame}

\begin{frame}{Task 3 - Q-learning (4 points)}
\small
Q(s,a)=1, reward=2, max next Q=3, gamma=0.5, and alpha=0.25.

\vspace{0.2cm}
\begin{enumerate}
\item[\textbf{a.}] Compute the nonterminal target and TD error. \hfill{\color{courseblue}\textbf{[2 points]}}
\item[\textbf{b.}] Compute the updated Q value and state the terminal change. \hfill{\color{courseblue}\textbf{[2 points]}}
\end{enumerate}
\end{frame}

\begin{frame}{Task 4 - Planning versus learning (4 points)}
\small
Compare value iteration and Q-learning for the same gridworld.

\vspace{0.2cm}
\begin{enumerate}
\item[\textbf{a.}] What information does each method require? \hfill{\color{courseblue}\textbf{[2 points]}}
\item[\textbf{b.}] Why must a Q-learning evaluation run use epsilon zero? \hfill{\color{courseblue}\textbf{[2 points]}}
\end{enumerate}
\end{frame}

\begin{frame}{Edge cases and defensive programming}
\small
\begin{itemize}
  \item A minimax root is terminal.
  \item An MDP action has probabilities that do not sum to one.
  \item A Q-learning transition is terminal.
\end{itemize}
\end{frame}

\begin{frame}{Submission instructions}
\small
\begin{itemize}
  \item Show every backup and update.
  \item State any convention needed for terminal reward.
  \item Submit only the completed response PDF.
\end{itemize}
\vspace{0.2cm}
\alert{Do not submit credentials, generated caches, or unapproved libraries.}
\end{frame}

\begin{frame}{Grading rubric}
\scriptsize
\begin{tabularx}{\textwidth}{>{\bfseries}p{0.28\textwidth} p{0.08\textwidth} X}
\toprule
Category & Points & Required evidence \\
\midrule
Minimax and alpha-beta & 6 & Tree trace. \\
MDP backup & 6 & Expected-return calculation. \\
Q-learning & 4 & One numerical update. \\
Planning versus learning & 4 & Short explanation. \\
\bottomrule
\end{tabularx}
\end{frame}

\begin{frame}{Reflection prompt}
In one sentence:

State one reason a learned policy can differ from a planning oracle.
\end{frame}

\end{document}
